% ======================================================================
% SECTION 1: INTRODUCTION
% Location in paper: First body section after the Abstract. Three
% thematic blocks. The introduction must cite and further address
% current patient legislation affecting oncology trials per the project
% brief.
% ======================================================================

[PROCESSING INSTRUCTION (Introduction block 1 - The Patient Priority
Thesis, 4 to 6 paragraphs): Open with a focused statement of the
patient priority thesis: the cancer patient, NOT the doctor, the
nurse, the trial sponsor, the IRB, or the regulator, is the priority
participant in any oncology clinical trial; advanced AI and advanced
robotics now offer the patient the opportunity to take direct control
over discovery, scheduling, robot selection, procedural modification,
real-time data feedback, and self-custody of health data; this paper
introduces seven proposed U.S. federal bills (HR 4501 through HR 4507)
that would enable that control.

Read patients/paper/Deep-Research-3/part1\_legal\_baseline.md in full
and quote the operational definition of patient control across six
functional dimensions: trial discovery, enrollment comprehension,
location choice, data and biospecimen governance, toxicity warnings,
and AI-exclusion contestability. Add a seventh dimension introduced
in this paper: physical-AI procedure execution choice (the patient
selects which surgical robot, humanoid, or companion robot performs
their care, drawing from the 10 robot categories enumerated in
patients/patient\_robot\_instructions\_fixed.tex pages 1 through 10 and
in national-platform/patient\_robot/01\_preamble\_robots1to5.tex and
02\_robots6to10\_closing.tex).

Cite the approximate 7 percent participation-rate statistic from
Deep-Research-3/part1\_legal\_baseline.md and frame it as the central
operational deficit that the new bills address. Cite
\cite{hhs_45cfr46_protection_human_subjects},
\cite{fda_informed_consent_guidance_2023}, and
\cite{fda_dct_guidance_2024}.

Close this block with a respectful framing of the FDA. The patient
priority thesis does NOT critique the FDA; it builds opportunistically
on the FDA's own April 28 2026 Real-Time Clinical Trials announcement
at new-trial/national-24-7-trial/FDA-April-2026/FDA\_RealTime\_Clinical\_Trials.md
and the FDA's 2024 Decentralized Clinical Trials Final Guidance. The
new bills are an opportunistic response to FDA-validated openings,
not an adversarial proposal. Cite \cite{fda_rtct_press_2026} and
\cite{fda_dct_guidance_2024}.]

[PROCESSING INSTRUCTION (Introduction block 2 - Current Patient
Legislation Affecting Oncology Trials, 6 to 9 paragraphs): Read the
following files in full and synthesize the layered legal stack against
which the seven proposed bills will operate. Each paragraph should
name one prior law or regulation by accurate name and Public Law (PL)
number, summarize what it provides, and identify the specific patient-
control deficit that one of the new bills addresses.

Files to read (all from the patients/paper/ Deep-Research pools):
- patients/paper/Deep-Research-2/chunk\_01\_intro\_ranks1-4.md (Ranks 1-4:
  21st Century Cures Act PL 114-255; CLINICAL TREATMENT Act PL 116-260;
  ACA Section 2709 PL 111-148; FDORA / FDA Modernization Act 2.0 PL
  117-328 Section 3209)
- patients/paper/Deep-Research-2/chunk\_02\_ranks5-9.md (Ranks 5-9: FDAAA
  PL 110-85 Section 801; Right to Try Act PL 115-176; FDARA PL 115-52;
  CA SB 37 (2001); ONC HTI-1 Final Rule 2023)
- patients/paper/Deep-Research-2/chunk\_03\_ranks10-13\_future.md
  (Ranks 10-13: AI/CDS guidance documents; diversity action plans;
  AI evidence; state law)
- patients/paper/Deep-Research-3/part1\_legal\_baseline.md (the
  Common Rule, broad-consent permissions, Cures Act/HIPAA data-access,
  FDA decentralized-trial guidance)
- patients/paper/Deep-Research-3/part2\_ai\_robotics\_legislation.md
  (the seven legislative vehicle types ranked by patient-control
  impact)

Cover, in order, with one short paragraph per item:

(2.a) 21st Century Cures Act, PL 114-255 (2016): interoperability,
information-blocking, real-world evidence, and the non-device clinical
decision support software carve-out. Identify the gap that HR 4507
addresses: information-blocking penalties exist
(\cite{federal_register_info_blocking_disincentives_2024}) but
patient-initiated self-custody of trial-eligibility data is not yet
guaranteed. Cite \cite{rank1_cures_act_2016} and
\cite{onc_cures_act_final_rule_page}.

(2.b) CLINICAL TREATMENT Act, PL 116-260 (2020): Medicaid coverage of
routine costs in qualifying trials from January 1 2022. Identify the
gap that HR 4501 addresses: AI trial-matching tools (TrialGPT, PRISM,
TrialMatchAI) become actionable only when Medicaid-eligible patients
can also select and self-book. Cite
\cite{rank2_clinical_treatment_act_2020}, \cite{jin2024trialgpt},
\cite{gupta2024prism}, and \cite{abdallah2026trialmatchai}.

(2.c) Affordable Care Act Section 2709, PL 111-148 (2010): private
insurance baseline for routine clinical-trial cost coverage. Identify
the gap that HR 4501 also addresses: payment parity for the
decentralized and at-home components of trials remains unprotected
under ACA Section 2709. Cite \cite{rank3_aca_clinical_trials_2010}
and \cite{usc_300gg8_clinical_trials}.

(2.d) FDORA / FDA Modernization Act 2.0, PL 117-328 (2022),
Section 3209: in silico and computer-model alternatives to animal
testing; mandate on diversity action plans; mandate on FDA decentralized
trial guidance. Identify the gap that HR 4506 addresses: federal
authorization for sustained American leadership in medical AI and
robotics has not yet been enacted as a standalone statute. Cite
\cite{rank4_fdora_fda_modernization_2022} and \cite{fda_dct_guidance_2024}.

(2.e) FDAAA Section 801, PL 110-85 (2007): ClinicalTrials.gov
infrastructure. Identify the gap that HR 4501 and HR 4507 jointly
address: structured eligibility data for patient-initiated discovery
remains incomplete. Cite \cite{rank5_fdaaa_2007}.

(2.f) ONC HTI-1 Final Rule (2023) and HTI-1 DSI Fact Sheet:
algorithm-transparency requirements for predictive decision support
interventions in certified health IT. Identify the gap that HR 4504
addresses: transparency exists, but mandatory error-rate tracking and
patient-initiated AI-second-opinion rights are not yet codified. Cite
\cite{rank9_onc_hti1_2023} and \cite{hti1_dsi_fact_sheet_2023}.

(2.g) FDA Informed Consent Guidance 2023, FDA/OHRP Key Information
Draft 2024, and OHRP Broad Consent Recommendations 2017. Identify the
gap that HR 4503 addresses: consent and broad-consent frameworks exist
but do not include a positive procedural-modification right where the
patient may modify their own treatment under physical-AI execution.
Cite \cite{fda_informed_consent_guidance_2023},
\cite{fda_ohrp_key_information_draft_2024}, and
\cite{ohrp_broad_consent_attachment_c_2017}.

(2.h) HHS HIPAA Right-to-Access (2025), ASTP/ONC Patient Portals
Brief 2024, and ASTP/ONC Hospital APIs Brief 2024. Identify the gap
that HR 4507 addresses: access exists but patient-initiated self-
custody of trial-eligibility-relevant pathology and genomics data is
not yet guaranteed. Cite \cite{hhs_right_to_access_2025},
\cite{astp_onc_patient_portals_apps_2024}, and
\cite{astp_onc_hospital_apis_2024}.

(2.i) FDA Real-Time Clinical Trials Press Announcement (April 28 2026)
and FDA Digital Health Technologies for Drug Development (2026).
Identify the openings that HR 4505 builds on: the FDA has validated a
real-time signal-sharing conduit (Paradigm Health) and named two
proof-of-concept trials (TRAVERSE and STREAM-SCLC). HR 4505 extends
this opening into a patient-initiated direct sponsor communication
right. Cite \cite{fda_rtct_press_2026} and
\cite{fda_dhts_drug_development_2026}.]

[PROCESSING INSTRUCTION (Introduction block 3 - Transition to the
Seven Bills, 2 to 3 paragraphs): Compose a transition paragraph that
maps each prior-legislation gap identified in block 2 onto the
corresponding new bill. Use a short table with three columns
(Prior Law, Patient-Control Gap, Proposed New Bill) generated using
the L{w} column type from main.tex; left-justify text in every column
to avoid right-margin overflow.

The second paragraph should preview each of the seven new bills by
short title only, in the order they appear in the paper (HR 4501
through HR 4507). Each preview is one sentence and names the prior
legislation that the bill adapts or revises. The seventh bill
(HR 4506) is the American Leadership in Medical AI and Robotics Act,
which captures the United States Number 1 framing required by the
project brief.

Close the introduction by stating that each subsequent section
(Sections 3 through 9 of this paper) presents one bill in the
legislative-act layout used by the prior site documentation template
at new-trial/site/all-documents/all\_documents.tex (Findings,
Definitions, Authorization or Rights, Implementation, Reporting and
Enforcement). Note that for individual patients the bills define
specific rights; for broad adoption across all U.S. cancer patients
in clinical trials the bills define federal preemption rules and
implementation timelines. Cite \cite{kawchak_2026_19176370} for the
prior site template and \cite{kawchak_2026_20045457} for this paper's
Zenodo DOI.]
